\section{Ograniczenia przeprowadzonych badań}
\label{sec:OgraniczeniaBadan}

Przeprowadzone badania umożliwiły porównanie działania wybranych algorytmów 
sortowania równoległego w~określonych warunkach sprzętowych i~programowych. 
Otrzymane rezultaty należy jednak rozpatrywać z~uwzględnieniem ograniczeń związanych z~przyjętą metodyką badań.

Pierwszym ograniczeniem było zastosowanie jednego środowiska sprzętowego. 
Wszystkie eksperymenty przeprowadzono na tej samej platformie wyposażonej w~procesor Intel Core i9-9980XE posiadający
18 rdzeni fizycznych i~36 wątków logicznych oraz 128~GB pamięci RAM.
Oznacza to, że uzyskane rezultaty odzwierciedlają zachowanie badanych implementacji na konkretnej konfiguracji sprzętowej. 
W przypadku innych procesorów, cechujących się inną liczbą rdzeni, organizacją pamięci podręcznej lub przepustowością pamięci, 
uzyskane relacje pomiędzy analizowanymi wariantami nie mogą być bezpośrednio uogólniane na inne konfiguracje sprzętowe.

Badania przeprowadzono również w~jednym środowisku systemowym.
Wykorzystano system Windows, w~którym moduł \texttt{multiprocessing} uruchamia procesy z~wykorzystaniem mechanizmu \texttt{spawn}.
Uzyskane wyniki uwzględniają zatem narzut właściwy dla tego sposobu uruchamiania procesów i~nie należy ich bezpośrednio porównywać 
z~rezultatami możliwymi do uzyskania w~innych systemach operacyjnych lub przy zastosowaniu innych metod tworzenia procesów.

Kolejnym ograniczeniem był zakres badanych rozmiarów danych.
W badaniach wykorzystano cztery rozmiary zbiorów danych: od 1~000 do 1~000~000 elementów.
Z kolei warianty równoległe badano dla 2, 4, 8, 16 oraz 32 jednostek wykonawczych.
Dla części algorytmów największy badany zbiór mógł być nadal niewystarczający, aby w~pełni wykorzystać większą liczbę jednostek wykonawczych.
Natomiast w~przypadku najmniejszych zbiorów bardzo krótki czas wykonania ograniczał liczbę zarejestrowanych próbek wykorzystania CPU. 
W~efekcie w~niektórych pomiarach odnotowano wartość 0\%, co nie oznacza jednak, że procesor nie był w~tym czasie wykorzystywany.
Z~tego względu wyniki uzyskane dla skrajnych rozmiarów danych należy interpretować z~uwzględnieniem tych ograniczeń.

Na otrzymane rezultaty wpływ miały również zastosowane progi przejścia do~wykonania sekwencyjnego.
Wartości progów zostały ustalone przed przeprowadzeniem końcowej serii eksperymentów 
i~nie były później modyfikowane na podstawie uzyskiwanych rezultatów. 
Umożliwiło to zachowanie jednakowych warunków pomiarowych, jednak uzyskane wyniki pokazują, 
że przyjęte progi nie były jednakowo korzystne dla wszystkich kombinacji rozmiaru danych i~liczby jednostek wykonawczych. 
W~niektórych konfiguracjach ograniczały one faktyczny stopień równoległości.
Konfiguracje te charakteryzowały się również niższymi wartościami przyspieszenia i~efektywności.

Uzyskane wyniki należy również interpretować z~uwzględnieniem specyfiki zastosowanych implementacji. 
Badane algorytmy różniły się między innymi pod względem podziału danych, 
organizacji etapów równoległych oraz zarządzania procesami i~pamięcią współdzieloną. 
W rezultacie zaobserwowane różnice w~czasach wykonania wynikały zarówno z~charakterystyki poszczególnych algorytmów, 
jak i~zastosowanych w~nich rozwiązań implementacyjnych.

Istotnym ograniczeniem pozostaje również charakter środowiska wykonawczego języka Python. 
Zastosowanie modułu \texttt{multiprocessing} umożliwiało wykonywanie zadań w~niezależnych procesach, 
jednak rzeczywisty poziom równoległości zależał również od sposobu organizacji procesów, 
przyjętych progów oraz mechanizmu planowania zadań przez system operacyjny.
Przeprowadzona dodatkowa diagnostyka wykazała, że liczba zadeklarowanych jednostek wykonawczych 
nie odpowiadała liczbie faktycznie tworzonych procesów potomnych w~przypadku Quick Sort i~Merge Sort.